\subsection{QC-010: padded injectivity for small total tail weight}

\paragraph{Statement QC-010.}
For the pair $(m,n)=(8,9)$, let
$\lambda=(9d-w,\beta_1,\beta_2,\beta_3)$ be a partition with at
most four rows, where $w=|\beta|$.
Then the padded multiplicity equals the full ambient multiplicity,
$q_\lambda=a_\lambda$, in either of the following regions:
\[
  w\leq25,\qquad\hbox{or}\qquad w=26,\ \beta_1\leq11.
  \tag{QC-010.1}
\]
The assertion concerns the whole highest-weight multiplicity space.
It does not assert a determinant-side inequality. The separate
one-column theorem S-007 also gives $q_\lambda=a_\lambda$ whenever
$\beta_1\leq9$, without a bound on $w$.

\paragraph{Centered description of the padded locus.}
The quaternary octic permanent image is the whole octic space.
Thus the padded image closure for $(8,9)$ is
$V=\{LQ_8\}$. On $F_0=1,F_1=0$, write the ordinary coefficient
expansion of a nonic as
\[
 F=x_0^9+\sum_{j=2}^9x_0^{9-j}B_j,\qquad
 B_j=\binom9j C_j.
\]
Here the $C_j$ are the normalized centered symmetric tensors.
Choose $C_2,\ldots,C_8$ freely and a linear form $l$ freely, and set
\[
 C_9=B_9=l^9+
       \sum_{j=2}^8(-1)^j\binom9j l^{9-j}C_j.
 \tag{QC-010.2}
\]
Then $F(-l,v)=0$ exactly, so $F=(x_0+l)Q_8$.
This gives actual points of the padded variety with freely
assignable lower centered tensors.

As proved by the localization argument in OBS-005, a homogeneous
highest-weight polynomial of tail $\beta$ has a unique centered
expression
$\sum_k F_0^{d-k}I_k(C_2,\ldots,C_9)$.
Setting $F_0=1$ identifies it injectively with a polynomial
$I=\sum_k I_k$ of total slot weight $w$ and $GL_3$ highest
weight $\beta$. It therefore suffices to prove injectivity
of substitution (QC-010.2) on the indicated centered spaces.

\paragraph{A quadratic polarization lemma.}
For a three-dimensional vector space $W$,
\[
 \operatorname{Sym}^2(\operatorname{Sym}^9W)
    =\bigoplus_{i=0}^4 S_{(18-2i,2i)}W.
 \tag{QC-010.3}
\]
The substitution
\[
 \sigma:\operatorname{Sym}^2(\operatorname{Sym}^9W)
       \longrightarrow
       \operatorname{Sym}^{10}W\otimes\operatorname{Sym}^{8}W,
 \qquad
 Q\longmapsto Q(l^9,lC_8)
 \tag{QC-010.4}
\]
is injective, with the harmless choice of dual conventions
understood consistently on both sides.

Here is a direct verification of both facts. Jacobi--Trudi gives
\[
 s_{(18-k,k)}=h_{18-k}h_k-h_{19-k}h_{k-1}.
\]
Summing for $0\leq k\leq9$ telescopes to $h_9^2$.
In the multiplicity-one summand $S_{(18-k,k)}$ a highest
symbol in two letters is
\[
 (x_0y_1-x_1y_0)^k x_0^{9-k}y_0^{9-k}.
\]
Interchanging the letters acts by $(-1)^k$, so precisely the
even $k$ occur in the symmetric square, proving (QC-010.3).
The tensor product on the right of (QC-010.4) contains every
$S_{(18-k,k)}$ for $0\leq k\leq8$, once, by the two-row
Pieri identity. The map is nonzero on each even-$k$ summand:
set $l=e_0$ and $C_8=e_0^{8-k}e_1^k$ in the displayed
highest symbol. Its relevant coefficient is nonzero.
An equivariant nonzero map on an irreducible module is
injective. The summands have distinct highest weights,
so their direct sum is also mapped injectively.

Tensoring (QC-010.4) with a polynomial coefficient ring
preserves injectivity. This last observation is essential:
the quadratic coefficients below may depend on other
centered tensors, but this creates no cancellation in
the independent variable $C_8$.

\paragraph{Proof when $w\leq25$.}
Because $C_9$ has slot weight nine, every polynomial of
weight at most $25$ has $C_9$-degree at most two. Write
\[
 I=I_0(C_2,\ldots,C_8)+I_1(C_2,\ldots,C_8)[C_9]
       +A(C_2,\ldots,C_8)[C_9,C_9].
 \tag{QC-010.5}
\]
The quadratic coefficient $A$ has slot weight $w-18\leq7$,
so it is independent of $C_8$.

Suppose substitution (QC-010.2) makes $I$ zero. Extract
the part linear in $C_8$ and of degree ten in $l$.
Neither $I_0$ nor $I_1[C_9]$ contributes, because their
degrees in $l$ are at most zero and nine. Since $A$ is
independent of $C_8$, the only contribution comes from
using $l^9$ once and $9lC_8$ once in the two $C_9$
arguments. It is the nonzero scalar $18$ times
$\sigma(A)$. All other pairs of lower polar terms either
have no $C_8$ or have a different degree in $l$.
The injectivity of (QC-010.4), over the coefficient ring
in $C_2,\ldots,C_7$, gives $A=0$.

Now the degree-nine part in $l$ is $I_1[l^9]$.
Pure ninth powers span $\operatorname{Sym}^9W$:
expanding $(l_1e_1+l_2e_2+l_3e_3)^9$ and comparing
monomial coefficients proves this directly in
characteristic zero. Thus $I_1=0$, and then $I_0=0$.
This proves the first region of (QC-010.1), without
any bound on $\beta_1$.

\paragraph{Reduction of the boundary $w=26$, $\beta_1\leq11$.}
There are still no cubic terms in $C_9$. In the quadratic
part, the only summand of (QC-010.3) whose first row can
be at most $11$ after tensoring with polynomial
representations is $S_{(10,8)}$.
Indeed tensoring with a polynomial representation never
removes boxes from the original diagram: such a
representation embeds in a finite direct sum of tensor
powers of $W$, and successive one-box Pieri rules
only add boxes.

The quadratic coefficient has slot weight eight and
therefore splits uniquely as
\[
 A=A_{<8}+L(C_8),
\]
where $A_{<8}$ uses only $C_2,\ldots,C_7$ and $L$ is
constant-linear in $C_8$. By Pieri, the portion of
$S_{(10,8)}W\otimes\operatorname{Sym}^8W$ with first
row at most eleven consists of exactly the following
six multiplicity-one summands:
\[
 (10,8,8),\ (10,9,7),\ (10,10,6),\
 (11,8,7),\ (11,9,6),\ (11,10,5).
 \tag{QC-010.6}
\]
Extracting $C_8$-degree two and $l$-degree ten in
(QC-010.2) isolates a map on these six summands.
We now check this map explicitly.

Let the symbolic letters $a,b,c$ have degrees $9,9,8$.
Write $\Delta=[abc]$ for the full three-dimensional
bracket, and $E=[ab],F=[ac],G=[bc]$ for the brackets
using the first two coordinates. The six highest
symbols, in the order (QC-010.6), are
\[
\begin{array}{c|l}
(10,8,8)&\Delta^8a_0b_0\\
(10,9,7)&\Delta^7E(a_0G+b_0F)\\
(10,10,6)&\Delta^6E^2FG\\
(11,8,7)&\Delta^7Ea_0b_0c_0\\
(11,9,6)&\Delta^6E^2c_0(a_0G+b_0F)\\
(11,10,5)&\Delta^5E^3FGc_0.
\end{array}
\tag{QC-010.7}
\]
Each has the stated multidegree and weight, is symmetric
in the two degree-nine letters, and is nonzero. Brackets
and free highest coordinates show the highest-weight
property directly.

To compute the isolated map, set $l=e_0$ and $a=e_0$,
differentiate once in $b_0$, and symmetrize $b,c$.
The derivative represents the symmetric tensor product
$lC_8$ in the second nonic argument; its normalization
factor is nonzero. Put
\[
 D=b_1c_2-b_2c_1,\qquad G=b_0c_1-b_1c_0.
\]
Up to nonzero scalar factors, the six outputs are
\[
 D^8,\quad 0,\quad D^6b_1^2c_1^2,\quad -D^7G,\quad
 D^6b_1c_1(b_1c_0+b_0c_1),\quad -D^5b_1^2c_1^2G.
 \tag{QC-010.8}
\]
Thus the map is injective on five of the six irreducible
summands. If $I$ vanishes on the padded locus, all five
corresponding coefficients of $L$ are zero. Only
$(10,9,7)$ needs an additional argument.

\paragraph{The exceptional summand $(10,9,7)$.}
Set every lower tensor except $C_5,C_8$ to zero.
The coefficient $A_{<8}$ vanishes, since weight eight
cannot be formed from weight five alone. The linear
coefficient $I_1$ has weight $26-9=17$ and also vanishes:
the equation $5u+8v=17$ has no nonnegative integer
solution. The constant term $I_0$ has no dependence
on $l$.

The exceptional symbol in (QC-010.7) has a nonzero
coefficient of $l$-degree eight, $C_5$-degree two,
and $C_8$-degree one after padded substitution.
Take $l=e_1$, $C_5=e_0^5$ and
$C_8=\operatorname{sym}(e_1e_2^7)$.
The two nonic arguments then use the terms
$l^4C_5$, with nonzero binomial normalization.
In the desired contraction the full brackets must
use all seven third-coordinate slots of $C_8$.
The relevant symbolic expression becomes
\[
  2[a b]^8a_0b_0c_1c_2^7.
\]
The coefficient at
$a_0^5a_1^4b_0^5b_1^4c_1c_2^7$ is exactly
\[
               2\binom84=140\ne0.
 \tag{QC-010.9}
\]
The fixed factorial and binomial tensor normalizations
are also nonzero in characteristic zero.
No remaining term can cancel this positive $l$-degree:
$A_{<8}$ and $I_1$ vanished on the specialization and
$I_0$ is independent of $l$. Consequently the
exceptional coefficient of $L$ is zero.

We have now proved $L=0$ in every case. Extract the
term linear in $C_8$ and of $l$-degree ten once more;
the injectivity lemma gives $A_{<8}=0$.
The degree-nine pure-power argument then gives
$I_1=I_0=0$. This proves the boundary region
and completes the proof of (QC-010.1).

\paragraph{Scope and verification.}
This is a characteristic-zero polynomial proof, not
a sampled rank inference. Its all-copy conclusion
uses the centered-space identification and the fact
that the entire highest-weight restriction kernel
has been shown to be zero. The source scripts and
their ranks in QC-008 are not needed by this proof.
The six highest symbols, their symmetry and
multidegrees, all five nonzero outputs and the
exceptional zero in (QC-010.8), and the coefficient
$140$ in (QC-010.9) were independently read and
recomputed algebraically by both the root and
support branches during the executor outage.
These were textual whole-proof reviews, not
new file or code executions.

